\scanpage{61}\mspage{9}
$x$ pt de $I\mathbin{\wedge}_{{}_2E_0}({}_4E_0)\smallsetminus I$

On trouve ainsi
\[
\Delta=\underbrace{\bigl(\underbrace{\underbrace{I\mathbin{\wedge}_{\underline M_0}\underline M}_{\text{degré }16}
\smallsetminus\underbrace{I}_{\text{degré }4}}_{\text{degré }12}\bigr)/\{1,\sigma\}}_{\text{degré }2}
\tag*{(27)}
\]
où on a posé maintenant
\[
\left\{\begin{aligned}
\underline M&\simeq {}_4E_0\quad\text{schéma en \unsure{groupe|groupes} loc. isom.}\\
&\hspace{5em}\text{à }(\mathbb Z/4\mathbb Z)^2_S\\
\underline M_0&\simeq\underline M\otimes_{\mathbb Z/4\mathbb Z}\mathbb Z/2
\simeq2\underline M\simeq{}_2E_0=\underline V
\end{aligned}\right.
\tag*{(28)}
\]
Ainsi \inserted{(dans le cas $E=E_0$, disons)} les pts d'ordre 4 de $E_0$ se groupent par \inserted{6} paires d'éléments symétriques, qui forment la configuration octaédrale $\Delta$ dans $X$.

On trouve ainsi le diagramme (presque cartésien)
\diagram{\begin{tikzcd}[ampersand replacement=\&,column sep={2.5em,between origins},row sep=1.5em]
{}\&\underline I\arrow[r,phantom,"\subset" description]\& E\arrow[rrrr,"\text{étale}"]\arrow[d,"{\substack{\text{ramifié au}\\\deleted{\unsure{I}}\text{dessus des pts}\\\text{de }\underline I}}"] \&{}\&{}\&{}\& E_0\;(\simeq E/\underline M_0)\arrow[d,"{\substack{\text{ramifié au dessus}\\\text{des pts de }D_J\cup t(S)}}"]\\
{(29)}\&\underline I\arrow[r,hook]\& X\arrow[rrrr,"{\text{ramifié au}}","{\substack{\text{dessus des}\\\text{pts de }D_J}}"']\&{}\&{}\&{}\&\Sigma_J\\[-1.4em]
{}\&{}\&\cup\&{}\&{}\&{}\&\cup\\[-1.4em]
{}\&{}\&\Delta\arrow[rrr,dash,end anchor=center]\&{}\&{}\&{}\arrow[r,dashed,start anchor=center]\&D_J
\end{tikzcd}}
qui est cartésien au dessus de $\unsure{\Sigma|\overline\Sigma}_J^*=\overline\Sigma_J\smallsetminus D_J$
% archive sheet number: 60
